\documentclass[11pt,a4paper]{ctexart}
\usepackage[margin=2.2cm]{geometry}
\usepackage{xcolor}
\usepackage{enumitem}
\usepackage{amssymb}
\usepackage[normalem]{ulem}
\usepackage{titlesec}
\usepackage{hyperref}
\definecolor{navy}{HTML}{1A365D}
\definecolor{ink}{HTML}{1F2933}
\definecolor{muted}{HTML}{66788A}
\definecolor{soft}{HTML}{EEF4FA}
\definecolor{gold}{HTML}{8A5A12}
\hypersetup{colorlinks=true,linkcolor=navy,urlcolor=navy}
\setlist{nosep}
\titleformat{\section}{\Large\bfseries\color{navy}}{}{0pt}{}
\titleformat{\subsection}{\large\bfseries\color{navy}}{}{0pt}{}
\pagestyle{plain}

\title{\textbf{七下物主代词与情态动词专项}\\[4pt]
\large 人教版（新课纲版）七年级下册 · 名词性物主代词 / 情态动词}
\author{初中英语资料站}
\date{2026-11-12}
\begin{document}
\maketitle

\section*{使用说明}
\begin{itemize}[leftmargin=1.4em]
\item 本讲义对应人教版（新课纲版）七年级下册语法要求：名词性物主代词与情态动词。
\item 共 4 讲语法、6 级 48 题。
\item 在线版：\url{https://1901151458-prog.github.io/english-materials/g7b-pronouns-modals/}
\end{itemize}

\section{名词性物主代词}
\begin{center}
\begin{tabular}{p{4.2cm}p{4.2cm}p{4.0cm}}
\hline
\textbf{形容词性} & \textbf{名词性} & \textbf{含义}\\
\hline
my / our & mine / ours & 我的 / 我们的\\
your & yours & 你的 / 你们的\\
his / her / its & his / hers / its & 他的 / 她的 / 它的\\
their & theirs & 他们的\\
\hline
\end{tabular}
\end{center}
\textbf{关键：}名词性物主代词相当于名词，\textbf{不能用于名词前}。Is this \underline{your} pen? --- No, the blue pen is \underline{mine}.

\section{情态动词 can / could}
\begin{center}
\begin{tabular}{p{2.6cm}p{9.8cm}}
\hline
\textbf{功能} & \textbf{例句}\\
\hline
能力 & We \underline{can speak} a little English. / He \underline{could swim} when he was six.\\
许可 & You \underline{can play} basketball at the weekend. / We \underline{can't run} in the hallways.\\
请求 & \underline{Can you help} me? / \underline{Could you tell} him to call me back?（更委婉）\\
\hline
\end{tabular}
\end{center}

\section{情态动词 may}
\begin{itemize}[leftmargin=1.4em]
\item 表示\textbf{许可}（可以）：\underline{May} I borrow your dictionary? / You \underline{may go} now.
\item 表示\textbf{可能性}（也许；可能）：I \underline{may be} late. / That \underline{may be} true.
\end{itemize}
\textbf{may 与 can：}两者都可表示许可，但 may 比 can \textbf{更正式、语气更委婉}。

\section{must 与 have to}
\begin{center}
\begin{tabular}{p{3.4cm}p{9.0cm}}
\hline
\textbf{区别} & \textbf{要点}\\
\hline
语气来源 & must 表主观看法、语气强；have to 强调客观需要\\
否定式 & mustn't ＝ 不准；禁止；don't have to ＝ 不必\\
人称数 & must 无变化；have to 有变化（has to）\\
\hline
\end{tabular}
\end{center}


\newpage
\section{分级练习（48 题）}
\noindent 按难度分级排列，先做基础，再做成分与扩展，最后做综合题。

\subsection*{第 1 级 · 名词性物主代词 · 形式}
\noindent\textcolor{gray}{mine / yours / his / hers / ours / theirs。}\par\smallskip
\begin{enumerate}[leftmargin=*,itemsep=4pt]
\item I 的名词性物主代词是 \_\_\_。
\begin{enumerate}[label=\Alph*.,leftmargin=2em,itemsep=1pt,topsep=2pt]
\item mine
\item my
\item me
\item I
\end{enumerate}
\item you 的名词性物主代词是 \_\_\_。
\begin{enumerate}[label=\Alph*.,leftmargin=2em,itemsep=1pt,topsep=2pt]
\item yours
\item your
\item you
\item yourself
\end{enumerate}
\item he 的名词性物主代词是 \_\_\_。
\begin{enumerate}[label=\Alph*.,leftmargin=2em,itemsep=1pt,topsep=2pt]
\item his
\item him
\item he
\item hers
\end{enumerate}
\item she 的名词性物主代词是 \_\_\_。
\begin{enumerate}[label=\Alph*.,leftmargin=2em,itemsep=1pt,topsep=2pt]
\item hers
\item her
\item she
\item him
\end{enumerate}
\item we 的名词性物主代词是 \_\_\_。
\begin{enumerate}[label=\Alph*.,leftmargin=2em,itemsep=1pt,topsep=2pt]
\item ours
\item our
\item us
\item we
\end{enumerate}
\item they 的名词性物主代词是 \_\_\_。
\begin{enumerate}[label=\Alph*.,leftmargin=2em,itemsep=1pt,topsep=2pt]
\item theirs
\item their
\item them
\item they
\end{enumerate}
\item it 的名词性物主代词是 \_\_\_。
\begin{enumerate}[label=\Alph*.,leftmargin=2em,itemsep=1pt,topsep=2pt]
\item its
\item it
\item it’s
\item itself
\end{enumerate}
\item 名词性物主代词相当于 \_\_\_。
\begin{enumerate}[label=\Alph*.,leftmargin=2em,itemsep=1pt,topsep=2pt]
\item 名词
\item 形容词
\item 动词
\item 副词
\end{enumerate}
\end{enumerate}

\subsection*{第 2 级 · 物主代词 · 辨析}
\noindent\textcolor{gray}{形容词性 vs 名词性：能不能独立使用。}\par\smallskip
\begin{enumerate}[leftmargin=*,itemsep=4pt]
\item A: Is this \_\_\_ pen?（你的）
\begin{enumerate}[label=\Alph*.,leftmargin=2em,itemsep=1pt,topsep=2pt]
\item your
\item yours
\item you
\item yourself
\end{enumerate}
\item The blue pen is \_\_\_.（我的）
\begin{enumerate}[label=\Alph*.,leftmargin=2em,itemsep=1pt,topsep=2pt]
\item mine
\item my
\item me
\item I
\end{enumerate}
\item Is it \_\_\_?（你的）
\begin{enumerate}[label=\Alph*.,leftmargin=2em,itemsep=1pt,topsep=2pt]
\item yours
\item your
\item you
\item yourself
\end{enumerate}
\item No, it isn’t \_\_\_.（我的）
\begin{enumerate}[label=\Alph*.,leftmargin=2em,itemsep=1pt,topsep=2pt]
\item mine
\item my
\item me
\item I
\end{enumerate}
\item No, \_\_\_ is at home.（她的）
\begin{enumerate}[label=\Alph*.,leftmargin=2em,itemsep=1pt,topsep=2pt]
\item hers
\item her
\item she
\item him
\end{enumerate}
\item 名词性物主代词 \_\_\_ 用于名词前。
\begin{enumerate}[label=\Alph*.,leftmargin=2em,itemsep=1pt,topsep=2pt]
\item 不能
\item 可以
\item 必须
\item 有时可以
\end{enumerate}
\item This is \_\_\_ book.（我的）
\begin{enumerate}[label=\Alph*.,leftmargin=2em,itemsep=1pt,topsep=2pt]
\item my
\item mine
\item me
\item I
\end{enumerate}
\item This book is \_\_\_.（我的）
\begin{enumerate}[label=\Alph*.,leftmargin=2em,itemsep=1pt,topsep=2pt]
\item mine
\item my
\item me
\item I
\end{enumerate}
\end{enumerate}

\subsection*{第 3 级 · 情态动词 can / could}
\noindent\textcolor{gray}{能力、许可、请求。}\par\smallskip
\begin{enumerate}[leftmargin=*,itemsep=4pt]
\item We \_\_\_ speak a little English.（能力）
\begin{enumerate}[label=\Alph*.,leftmargin=2em,itemsep=1pt,topsep=2pt]
\item can
\item may
\item must
\item have to
\end{enumerate}
\item He \_\_\_ swim when he was only six.
\begin{enumerate}[label=\Alph*.,leftmargin=2em,itemsep=1pt,topsep=2pt]
\item could
\item can
\item may
\item must
\end{enumerate}
\item You \_\_\_ play basketball at the weekend.（许可）
\begin{enumerate}[label=\Alph*.,leftmargin=2em,itemsep=1pt,topsep=2pt]
\item can
\item must
\item have to
\item need
\end{enumerate}
\item \_\_\_ you help me with my Chinese?（请求）
\begin{enumerate}[label=\Alph*.,leftmargin=2em,itemsep=1pt,topsep=2pt]
\item Can
\item May
\item Must
\item Have to
\end{enumerate}
\item \_\_\_ you tell him to call me back?（更委婉）
\begin{enumerate}[label=\Alph*.,leftmargin=2em,itemsep=1pt,topsep=2pt]
\item Could
\item Can
\item May
\item Must
\end{enumerate}
\item We \_\_\_ run in the hallways.（不许）
\begin{enumerate}[label=\Alph*.,leftmargin=2em,itemsep=1pt,topsep=2pt]
\item can’t
\item mustn’t
\item don’t have to
\item needn’t
\end{enumerate}
\item Father said I \_\_\_ swim in the pool.
\begin{enumerate}[label=\Alph*.,leftmargin=2em,itemsep=1pt,topsep=2pt]
\item could
\item can
\item may
\item must
\end{enumerate}
\item I \_\_\_ play badminton, but I can’t play volleyball.
\begin{enumerate}[label=\Alph*.,leftmargin=2em,itemsep=1pt,topsep=2pt]
\item can
\item may
\item must
\item have to
\end{enumerate}
\end{enumerate}

\subsection*{第 4 级 · 情态动词 may}
\noindent\textcolor{gray}{许可「可以」、可能性「也许」。}\par\smallskip
\begin{enumerate}[leftmargin=*,itemsep=4pt]
\item \_\_\_ I borrow your dictionary?（许可）
\begin{enumerate}[label=\Alph*.,leftmargin=2em,itemsep=1pt,topsep=2pt]
\item May
\item Must
\item Have to
\item Need
\end{enumerate}
\item Class is over. You \_\_\_ go now.（许可）
\begin{enumerate}[label=\Alph*.,leftmargin=2em,itemsep=1pt,topsep=2pt]
\item may
\item must
\item have to
\item need
\end{enumerate}
\item I \_\_\_ be late, so don’t wait for me.（也许）
\begin{enumerate}[label=\Alph*.,leftmargin=2em,itemsep=1pt,topsep=2pt]
\item may
\item must
\item have to
\item can
\end{enumerate}
\item That \_\_\_ be true.（可能）
\begin{enumerate}[label=\Alph*.,leftmargin=2em,itemsep=1pt,topsep=2pt]
\item may
\item must
\item can
\item need
\end{enumerate}
\item may 表示许可时比 can \_\_\_。
\begin{enumerate}[label=\Alph*.,leftmargin=2em,itemsep=1pt,topsep=2pt]
\item 更正式、更委婉
\item 更随便
\item 更强烈
\item 意思完全不同
\end{enumerate}
\item may 的两个主要意思：许可和 \_\_\_。
\begin{enumerate}[label=\Alph*.,leftmargin=2em,itemsep=1pt,topsep=2pt]
\item 可能性
\item 能力
\item 义务
\item 请求
\end{enumerate}
\item You \_\_\_ go now if you finish.
\begin{enumerate}[label=\Alph*.,leftmargin=2em,itemsep=1pt,topsep=2pt]
\item may
\item must
\item have to
\item need
\end{enumerate}
\item It \_\_\_ rain tomorrow.（也许）
\begin{enumerate}[label=\Alph*.,leftmargin=2em,itemsep=1pt,topsep=2pt]
\item may
\item must
\item can
\item need
\end{enumerate}
\end{enumerate}

\subsection*{第 5 级 · must 与 have to}
\noindent\textcolor{gray}{主观 vs 客观；mustn’t vs don’t have to。}\par\smallskip
\begin{enumerate}[leftmargin=*,itemsep=4pt]
\item We \_\_\_ be on time. Don’t arrive late.（主观）
\begin{enumerate}[label=\Alph*.,leftmargin=2em,itemsep=1pt,topsep=2pt]
\item must
\item have to
\item may
\item can
\end{enumerate}
\item We \_\_\_ be quiet in the library.（客观需要）
\begin{enumerate}[label=\Alph*.,leftmargin=2em,itemsep=1pt,topsep=2pt]
\item have to
\item must
\item may
\item can
\end{enumerate}
\item You \_\_\_ talk to your mother like that.（不准）
\begin{enumerate}[label=\Alph*.,leftmargin=2em,itemsep=1pt,topsep=2pt]
\item mustn’t
\item don’t have to
\item needn’t
\item can’t
\end{enumerate}
\item You \_\_\_ come if you don’t want to.（不必）
\begin{enumerate}[label=\Alph*.,leftmargin=2em,itemsep=1pt,topsep=2pt]
\item don’t have to
\item mustn’t
\item needn’t
\item can’t
\end{enumerate}
\item She \_\_\_ work late today.（第三人称）
\begin{enumerate}[label=\Alph*.,leftmargin=2em,itemsep=1pt,topsep=2pt]
\item has to
\item have to
\item must
\item may
\end{enumerate}
\item must 有 \_\_\_ 变化。
\begin{enumerate}[label=\Alph*.,leftmargin=2em,itemsep=1pt,topsep=2pt]
\item 没有人称和数
\item 有人称和数
\item 有过去式
\item 有复数
\end{enumerate}
\item mustn’t 的意思是 \_\_\_。
\begin{enumerate}[label=\Alph*.,leftmargin=2em,itemsep=1pt,topsep=2pt]
\item 不准；禁止
\item 不必
\item 不应该
\item 不能
\end{enumerate}
\item They \_\_\_ get to the station before five o’clock.
\begin{enumerate}[label=\Alph*.,leftmargin=2em,itemsep=1pt,topsep=2pt]
\item have to
\item has to
\item must has
\item have
\end{enumerate}
\end{enumerate}

\subsection*{第 6 级 · 综合提升}
\noindent\textcolor{gray}{物主代词与情态动词混合辨析。}\par\smallskip
\begin{enumerate}[leftmargin=*,itemsep=4pt]
\item This is \_\_\_ book.（我的）
\begin{enumerate}[label=\Alph*.,leftmargin=2em,itemsep=1pt,topsep=2pt]
\item my
\item mine
\item me
\item I
\end{enumerate}
\item That book is \_\_\_.（他的）
\begin{enumerate}[label=\Alph*.,leftmargin=2em,itemsep=1pt,topsep=2pt]
\item his
\item him
\item he
\item hers
\end{enumerate}
\item You \_\_\_ finish your homework first.（必须）
\begin{enumerate}[label=\Alph*.,leftmargin=2em,itemsep=1pt,topsep=2pt]
\item must
\item may
\item can
\item need
\end{enumerate}
\item \_\_\_ I use your phone?（委婉许可）
\begin{enumerate}[label=\Alph*.,leftmargin=2em,itemsep=1pt,topsep=2pt]
\item May
\item Must
\item Have to
\item Can’t
\end{enumerate}
\item We \_\_\_ hurry. The train is leaving in five minutes.（客观）
\begin{enumerate}[label=\Alph*.,leftmargin=2em,itemsep=1pt,topsep=2pt]
\item have to
\item must
\item may
\item can
\end{enumerate}
\item The pen isn’t \_\_\_. It’s \_\_\_ .（我的；她的）
\begin{enumerate}[label=\Alph*.,leftmargin=2em,itemsep=1pt,topsep=2pt]
\item mine; hers
\item my; her
\item mine; her
\item my; hers
\end{enumerate}
\item He \_\_\_ speak three languages.
\begin{enumerate}[label=\Alph*.,leftmargin=2em,itemsep=1pt,topsep=2pt]
\item can
\item may
\item must
\item has to
\end{enumerate}
\item You \_\_\_ be late for class.
\begin{enumerate}[label=\Alph*.,leftmargin=2em,itemsep=1pt,topsep=2pt]
\item mustn’t
\item don’t have to
\item needn’t
\item may not
\end{enumerate}
\end{enumerate}

\newpage
\section*{答案与解析}
\small
\begin{enumerate}[leftmargin=1.6em,itemsep=2pt]
\item \textbf{O1-1 A}：教材表2：my（形容词性）→ mine（名词性）。
\item \textbf{O1-2 A}：your → yours；教材表2。
\item \textbf{O1-3 A}：his 形容词性与名词性同形；教材表2。
\item \textbf{O1-4 A}：her → hers；教材表2。
\item \textbf{O1-5 A}：our → ours；教材表2。
\item \textbf{O1-6 A}：their → theirs；教材表2。
\item \textbf{O1-7 A}：its 形容词性与名词性同形；教材表2。
\item \textbf{O1-8 A}：教材：名词性物主代词相当于名词。
\item \textbf{O2-1 A}：修饰名词 pen 用形容词性 your；教材例句。
\item \textbf{O2-2 A}：独立使用作表语用名词性 mine；教材例句。
\item \textbf{O2-3 A}：独立使用用名词性 yours；教材例句。
\item \textbf{O2-4 A}：名词性 mine 独立使用；教材例句。
\item \textbf{O2-5 A}：名词性 hers 独立作主语；教材例句。
\item \textbf{O2-6 A}：教材：名词性物主代词不能用于名词前。
\item \textbf{O2-7 A}：名词前用形容词性 my。
\item \textbf{O2-8 A}：独立使用用名词性 mine。
\item \textbf{O3-1 A}：can 表示能力；教材例句。
\item \textbf{O3-2 A}：过去的能力用 could；教材例句。
\item \textbf{O3-3 A}：can 表示许可；教材例句。
\item \textbf{O3-4 A}：can 表示请求；教材例句。
\item \textbf{O3-5 A}：could 表示请求比 can 更委婉；教材例句。
\item \textbf{O3-6 A}：教材例句 We can’t run in the hallways 表许可的否定。
\item \textbf{O3-7 A}：主句过去时，用 could；教材例句。
\item \textbf{O3-8 A}：can 表能力；教材例句。
\item \textbf{O4-1 A}：may 表许可「可以」；教材例句。
\item \textbf{O4-2 A}：may 表许可；教材例句。
\item \textbf{O4-3 A}：may 表可能性「也许」；教材例句。
\item \textbf{O4-4 A}：may 表可能性；教材例句。
\item \textbf{O4-5 A}：教材：may 的用法比 can 更正式，语气更委婉。
\item \textbf{O4-6 A}：教材：may 表示许可（可以）和可能性（也许；可能）。
\item \textbf{O4-7 A}：may 表许可。
\item \textbf{O4-8 A}：may 表可能性。
\item \textbf{O5-1 A}：must 表主观看法、语气强；教材例句。
\item \textbf{O5-2 A}：have to 强调客观需要；教材例句。
\item \textbf{O5-3 A}：mustn’t 表示不准；禁止；教材原句。
\item \textbf{O5-4 A}：don’t have to 表示不必；教材原句。
\item \textbf{O5-5 A}：have to 有人称和数的变化，第三人称用 has to；教材例句。
\item \textbf{O5-6 A}：教材：must 没有人称和数的变化。
\item \textbf{O5-7 A}：教材：mustn’t 表示「不准；禁止」。don’t have to 才表示「不必」。
\item \textbf{O5-8 A}：主语 they 用 have to；教材例句。
\item \textbf{O6-1 A}：名词前用形容词性 my。
\item \textbf{O6-2 A}：his 形容词性与名词性同形。
\item \textbf{O6-3 A}：must 表示必须。
\item \textbf{O6-4 A}：May I ...? 委婉请求许可。
\item \textbf{O6-5 A}：强调客观需要，用 have to；教材例句。
\item \textbf{O6-6 A}：两空都独立使用，用名词性物主代词。
\item \textbf{O6-7 A}：can 表示能力。
\item \textbf{O6-8 A}：mustn’t 表示不准；禁止。
\end{enumerate}

\end{document}
